\section{Authorization is separate from declared identity}\label{sec:C19:authorization}
MCP's HTTP authorization specification distinguishes client, protected resource server, resource owner, and authorization server roles. Authorization is optional at the protocol-implementation level; HTTP implementations that support it SHOULD conform to the cited framework, while stdio implementations SHOULD instead obtain credentials from their environment~\cite{mcp2026auth}. A server being reachable does not imply that every capability is authorized, and a client being authenticated does not imply that every requested effect is within the user's task scope.

The current specification uses protected-resource and authorization-server discovery mechanisms and defined registration options~\cite{mcp2026auth}. These are protocol and credential contracts, not a complete application policy. An access token can permit a broad API scope while the current user request authorizes only a narrow task. A local monitor may need to enforce the narrower boundary even when the remote service accepts the credential.

The protocol's client and server information fields are self-reported metadata rather than cryptographic proof of identity~\cite{mcp2026discover}. Applications should not treat a familiar display name or version string as authorization evidence. Trust should derive from the configured endpoint, transport authentication, credential validation, and applicable policy. This is particularly important for discovery catalogs that combine many servers with similar tool names.

Delegation adds another question: whose authority is being exercised? A host can authenticate to a tool server using its own credential while acting on behalf of a user. A remote agent can request a capability that the original user never intended to delegate. The application should preserve the principal and scope through each hop, avoiding silent expansion when a worker or skill is loaded. Protocol compatibility makes the message exchange possible; it does not decide the appropriate delegation policy.

Credential exposure should not be treated as ordinary memory. The model generally needs a capability reference and allowed scope, not a plaintext token in its context. If a token or secret enters a tool result, subsequent summaries and indexes can spread it beyond its intended lifetime. A secure wrapper can keep credentials outside the model-visible evidence packet while recording which credential scope admitted an effect. The exact storage and redaction mechanism remains application-specific.

A proposed authorization test matrix includes correct credentials with excessive task scope, stale authorization after a user correction, catalog changes under different principals, and a reconnect that interleaves unrelated users. Measure both unauthorized acceptance and false rejection of legitimate work. A system that blocks every call is not a useful agent, while one that completes every task with broad credentials has not demonstrated least-privilege operation.

